\section{Related Work}

\paragraph{Low-rank tensor models.}
Classical tensor decompositions such as CP, Tucker, tensor train, and tensor
ring provide compact representations for multiway data
\citep{Kolda2009Tensor,DeLathauwer2000MLSVD,oseledets2011,zhao2016tr}. Tucker
is the closest baseline for our work because its core tensor and mode factors
give a flexible multilinear backbone. Most methods in this family improve the
representation by changing ranks, factor structures, regularizers, or
optimization procedures. NTD-PL keeps the same Tucker parameterization and
instead changes the observation map from the latent tensor to the measured
entries.

\paragraph{Nonlinear tensor factorization.}
Prior work has introduced nonlinearity through generalized likelihoods,
kernels, Gaussian processes, neural decoders, and nonlinear tensor network
architectures
\citep{hong2020generalized,xu2011infinite,zhe2016distributed,pan2020streaming,saragadam2024deeptensor,varolgunes2023nmtucker}.
These approaches are often more expressive than polynomial links, but they
can trade off structural transparency, optimization simplicity, or explicit
finite-dimensional analysis. Our goal is complementary. We use a minimal
shared observation link so that the nonlinear effect can still be analyzed as
a structured extension of Tucker rather than as an unconstrained decoder.

\paragraph{Polynomial feature lifts and high-order interactions.}
Polynomial expansions and Volterra-style models are classical tools for
representing nonlinear systems. In tensor settings, high-order interactions
can be represented explicitly but often incur rapidly growing parameter and
contraction costs. Polynomial matrix and tensor factorization methods use
feature expansions to increase expressivity, but they typically introduce
larger lifted spaces, independent high-order parameters, or model-specific
factor structures. NTD-PL can be viewed as a shared-backbone compression of
high-order Tucker interactions: the lifted features are induced by powers of
one learned latent tensor, not by an independent high-order tensor parameter.

\paragraph{Rank lifting and post-hoc calibration.}
Two natural alternatives explain why the empirical comparisons are organized
as mechanism checks. Rank lifting gives Tucker more ordinary multilinear
capacity and can absorb part of the nonlinear residual, but it enlarges the
core and factors throughout the model. Post-hoc scalar calibration is cheaper,
but it remaps a Tucker reconstruction after the latent geometry has already
been learned. NTD-PL sits between these alternatives: it adds only a small
shared link, but learns that link jointly with the Tucker factors.
